\chapter{Transformer 核心公式彻底推导：每一个符号从何而来}

\begin{mathBox}[核心数学定理推导]
\textbf{阅读前须知}：在这篇教程中，我们将直面大模型皇冠上的明珠——Transformer 的数学核心。我们不会像很多教材那样直接甩给你公式，而是像侦探破案一样，带你一步一步把每一个符号、每一个分母、每一个负号的数学来源推导得清清楚楚。
\end{mathBox}

\vspace{0.8em}\hrule\vspace{0.8em}

\section{[架构] 一、终极核心公式全貌}

在所有大模型（GPT-4、Llama-3、Qwen-2.5）的心脏位置，跳动着的都是这行经典的注意力公式：

\[
\text{Attention}(Q, K, V) = \text{softmax}\left(\frac{Q K^T}{\sqrt{d_k}}\right) V
\]

请盯住这个公式 3 秒钟。它包含 5 个核心要素：
\begin{enumerate}
  \item \textbf{$Q$（Query，查询）}
  \item \textbf{$K$（Key，键）}
  \item \textbf{$V$（Value，值）}
  \item \textbf{$Q K^T$（点积打分矩阵）}
  \item \textbf{分母 $\sqrt{d_k}$（缩放因子，最精妙的数学设计）}
\end{enumerate}

下面我们逐一拆解。

\vspace{0.8em}\hrule\vspace{0.8em}

\section{[观察] 二、为什么需要拆分成 Q、K、V？（搜索引擎直觉）}

很多人初学时都有一个巨大的困惑：\textbf{明明输入的只有一句话的词嵌入 $X$，为什么偏偏要用三个不同的矩阵把 $X$ 变成 $Q$、$K$、$V$？直接拿 $X$ 算不行吗？}

\subsection{类比：YouTube / 百度搜索引擎的底层逻辑}

想象你在 YouTube 搜索框里搜视频：
\begin{enumerate}
  \item \textbf{你的搜索词}：\texttt{"如何学习大模型推理"} —— 这就是 \textbf{Query（查询）}。它代表当前词“迫切想寻找什么信息”。
  \item \textbf{库里每个视频的标题/标签}：\texttt{"大模型实战教程"}, \texttt{"美食烹饪视频"} —— 这就是 \textbf{Key（索引键）}。它代表每个词用来被别人检索的“特征标签”。
  \item \textbf{视频真正的画面和声音}：几十分钟的真实视频内容 —— 这就是 \textbf{Value（真实值）}。它代表每个词真正承载的信息。
\end{enumerate}

\begin{figure}[htbp]
\centering
\begin{tikzpicture}[scale=0.88, >=stealth,
  box/.style={rectangle, rounded corners=3pt, minimum width=2.6cm, minimum height=0.75cm, text centered, font=\sffamily\footnotesize\bfseries, draw=primary, fill=primary!8}
]
  \node[box, fill=secondary!15, draw=secondary] (X) at (0, 0) {输入向量序列 $X$};
  
  \node[box] (Q) at (4.5, 1.6) {Query ($Q = X W_q$)};
  \node[box] (K) at (4.5, 0) {Key ($K = X W_k$)};
  \node[box] (V) at (4.5, -1.6) {Value ($V = X W_v$)};
  
  \node[anchor=west, font=\sffamily\scriptsize, color=darktext] at (6.2, 1.6) {当前词想要探寻什么信息？};
  \node[anchor=west, font=\sffamily\scriptsize, color=darktext] at (6.2, 0) {当前词有什么特征可被检索？};
  \node[anchor=west, font=\sffamily\scriptsize, color=darktext] at (6.2, -1.6) {当前词承载的实际内容是什么？};
  
  \draw[->, thick, primary] (X.east) -- ++(1.2, 0) |- (Q.west);
  \draw[->, thick, primary] (X.east) -- (K.west);
  \draw[->, thick, primary] (X.east) -- ++(1.2, 0) |- (V.west);
\end{tikzpicture}
\caption{自注意力机制三元组 $Q, K, V$ 投影与语义物理意义}
\end{figure}

\subsection{为什么不能直接用 $X \cdot X^T$？}

如果直接拿原始特征自己和自己做点积：
\begin{itemize}
  \item 相似度矩阵就变成了对称矩阵：词 A 对词 B 的注意力，必须强制等于词 B 对词 A 的注意力。
  \item 但在人类语言中，\textbf{关联往往是极其不对称的}！
  \item 例如句子：“小猫跳上了桌子，因为它饿了”。
  \item 当模型读到代词“它”时，“它”需要极其强烈地关注前文的“小猫”（消歧代词）；
  \item 但前面的“小猫”在被输入时，根本不需要关注后面很远的“它”。
  \item \textbf{引入独立的 $W_q$ 和 $W_k$ 投影，彻底打破了对称性束缚，让信息流动的指向性更加灵活、自由！}
\end{itemize}

\vspace{0.8em}\hrule\vspace{0.8em}

\section{[爆炸] 三、【重磅数学推导】：为什么必须除以 $\sqrt{d_k}$？}

这是面试中最高频、最考验底层功底的一道经典推导题：\textbf{论文作者为什么不除以 $d_k$，也不除以 10，偏偏要除以根号 $\sqrt{d_k}$？}

我们来进行完整的\textbf{方差演化数学推导}。

\subsection{假设输入是标准正态分布}

设 Query 向量的一个分量为 $q_i$，Key 向量的一个分量为 $k_i$。
在经过层归一化（LayerNorm）后，模型内部特征通常满足均值为 0、方差为 1 的独立标准正态分布：
\begin{itemize}
  \item 期望（均值）：$E[q_i] = 0, \quad E[k_i] = 0$
  \item 方差：$\text{Var}(q_i) = 1, \quad \text{Var}(k_i) = 1$
  \item 单个维度的长度为 $d_k$（例如 $d_k = 64$ 或 $128$）。
\end{itemize}

\subsection{计算点积 $Z = \vec{q} \cdot \vec{k}$ 的均值与方差}

点积公式展开为：

\[
Z = \sum_{i=1}^{d_k} q_i k_i = q_1 k_1 + q_2 k_2 + \dots + q_{d_k} k_{d_k}
\]

\subsubsection{(1) 计算点积的数学期望 $E[Z]$：}

根据独立变量期望的可加性与可乘性：

\[
E[Z] = \sum_{i=1}^{d_k} E[q_i k_i] = \sum_{i=1}^{d_k} E[q_i] \cdot E[k_i] = \sum_{i=1}^{d_k} (0 \times 0) = 0
\]

均值仍然是 0，没有偏置。

\subsubsection{(2) 计算单项乘积 $q_i k_i$ 的方差 $\text{Var}(q_i k_i)$：}

根据方差定义公式 $\text{Var}(X) = E[X^2] - (E[X])^2$：

\[
\text{Var}(q_i k_i) = E[(q_i k_i)^2] - (E[q_i k_i])^2 = E[q_i^2] E[k_i^2] - 0
\]

因为 $E[q_i^2] = \text{Var}(q_i) + (E[q_i])^2 = 1 + 0 = 1$，同理 $E[k_i^2] = 1$：

\[
\text{Var}(q_i k_i) = 1 \times 1 = 1
\]

\subsubsection{(3) 计算 $d_k$ 项累加后的总方差 $\text{Var}(Z)$：}

由于这 $d_k$ 个维度相互独立，总方差等于各项方差相加：

\[
\text{Var}(Z) = \sum_{i=1}^{d_k} \text{Var}(q_i k_i) = \underbrace{1 + 1 + \dots + 1}_{d_k \text{ 个}} = d_k
\]

\textbf{根据标准差是方差的平方根}：

\[
\sigma(Z) = \sqrt{\text{Var}(Z)} = \sqrt{d_k}
\]

\vspace{0.8em}\hrule\vspace{0.8em}

\subsection{方差爆炸的灾难后果：Softmax 梯度消失}

请看当隐藏维度变大时发生的事情：
\begin{itemize}
  \item 现代大模型每个头的维度 $d_k$ 通常是 \textbf{$128$}。
  \item 点积未缩放前的标准差就是 $\sqrt{128} \approx \mathbf{11.31}$！
  \item 按照正态分布 $3\sigma$ 原则，大部分点积数值会在 $[-34, +34]$ 之间剧烈震荡。
\end{itemize}

如果把 $[+30, -10, 0]$ 送入 Softmax 会发生什么？
\begin{itemize}
  \item $e^{30} \approx 1.06 \times 10^{13}$
  \item $e^{0} = 1$
  \item 最大的一项占据了 $99.99999999999\%$ 的概率，其他所有项几乎全被压为 $0$！
\end{itemize}

\textbf{更致命的是梯度反向传播公式}：
Softmax 的导数公式为：

\[
\frac{\partial p_i}{\partial z_j} = p_i (\delta_{ij} - p_j)
\]

当某个 $p_i \approx 1$ 而其他 $p_j \approx 0$ 时，这个导数项\textbf{直接逼近于 0}！
\textbf{梯度彻底消失，网络所有参数锁死，无法学习任何深层知识！}

\vspace{0.8em}\hrule\vspace{0.8em}

\subsection{救世主：除以 $\sqrt{d_k}$}

根据随机变量的方差性质：常数 $c$ 提出来要平方，即 $\text{Var}(c X) = c^2 \text{Var}(X)$。
如果我们将点积 $Z$ 除以 $\sqrt{d_k}$：

\[
\text{Var}\left(\frac{Z}{\sqrt{d_k}}\right) = \left(\frac{1}{\sqrt{d_k}}\right)^2 \cdot \text{Var}(Z) = \frac{1}{d_k} \cdot d_k = \mathbf{1}
\]

\textbf{完美的代数闭环！}
无论模型的维度 $d_k$ 做得多大（64、128、256），除以 $\sqrt{d_k}$ 之后，点积结果的方差被\textbf{死死压制在 1.0}！
使得输入进 Softmax 的数值始终落在 $[-3, +3]$ 的黄金活跃区间，Softmax 既能分出高下，又保持了充沛的梯度流动！

\vspace{0.8em}\hrule\vspace{0.8em}

\section{[禁用] 四、因果掩码（Causal Mask）：为什么必须填充 $-\infty$？}

在生成文本时，第 1 个词绝对不能看到第 2 个词的内容。
我们通过一个掩码矩阵 $M$ 与原始点积得分相加：

\[
S_{\text{masked}} = \frac{Q K^T}{\sqrt{d_k}} + M
\]

其中上三角区域（代表未来时刻）全部被填入 $-\infty$（代码中通常用 $-10000.0$ 或 \texttt{-1e9}）：

\[
M = \begin{bmatrix}
0 & -\infty & -\infty \\
0 & 0 & -\infty \\
0 & 0 & 0
\end{bmatrix}
\]

\subsection{数学逻辑：}

根据极限：

\[
\lim_{x \to -\infty} e^x = 0
\]

在计算第 1 行的 Softmax 时：

\[
p_{1, 2} = \frac{e^{S_{1, 2} + (-\infty)}}{\sum e^{\dots}} = \frac{0}{\sum e^{\dots}} = \mathbf{0}
\]

\textbf{加法叠加负无穷，经指数变换后精确转化为乘法意义上的 0 权重，彻底斩断对未来信息的窥探！}

\vspace{0.8em}\hrule\vspace{0.8em}

\section{[路径] 五、残差连接（Residual Connection）：信息的高速公路}

在 Transformer 的每一层后面，都有这样一行代码：

\[
X_{\text{output}} = X_{\text{input}} + \text{SubLayer}(X_{\text{input}})
\]

\subsection{为什么不能直接 $X_{\text{output}} = \text{SubLayer}(X_{\text{input}})$？}

假设一个大模型有 80 层（如 Llama-3-70B）：
在反向传播计算梯度时，根据微积分链式法则：

\[
\frac{\partial \mathcal{L}}{\partial X_0} = \frac{\partial \mathcal{L}}{\partial X_{80}} \times \prod_{l=1}^{80} \frac{\partial X_l}{\partial X_{l-1}}
\]

80 个矩阵导数连乘，只要每一层的雅可比矩阵范数稍微小于 1（例如 0.9），$0.9^{80} \approx 0.0002$，底层的参数梯度瞬间衰减归零，大模型根本训练不起来！

\subsection{加入残差后的神奇变化：}

因为 $X_l = X_{l-1} + F(X_{l-1})$，对 $X_{l-1}$ 求导：

\[
\frac{\partial X_l}{\partial X_{l-1}} = \mathbf{I} + \frac{\partial F(X_{l-1})}{\partial X_{l-1}}
\]

展开到前向损失：

\[
\frac{\partial \mathcal{L}}{\partial X_{l-1}} = \frac{\partial \mathcal{L}}{\partial X_l} \left( \mathbf{I} + \frac{\partial F}{\partial X_{l-1}} \right) = \frac{\partial \mathcal{L}}{\partial X_l} + \frac{\partial \mathcal{L}}{\partial X_l} \frac{\partial F}{\partial X_{l-1}}
\]

看到了吗？式子中永远出现了一个干净纯粹的单位矩阵 $\mathbf{I}$！
\textbf{它宛如一条畅通无阻的高速公路，无论模型叠到 100 层还是 1000 层，梯度都能沿着这条高速路无损地一路冲回第一层！}

\vspace{0.8em}\hrule\vspace{0.8em}

\section{[循环] 六、RoPE（旋转位置编码）：复平面上的几何芭蕾}

最后我们来看现代开源模型（LLaMA/Qwen/DeepSeek）一致采用的 \textbf{RoPE（Rotary Position Embedding）}。

\subsection{核心目标}

我们希望找到一个变换函数 $R(x, m)$，给向量 $x$ 注入位置序号 $m$。
并且满足一个梦幻般的特性：\textbf{两个向量注入位置 $m$ 和 $n$ 后的点积，只和它们的相对距离 $(m - n)$ 有关！}

\[
\langle R(q, m), R(k, n) \rangle = g(q, k, m - n)
\]

\subsection{二维复平面的几何推导}

在复数几何中，将一个二维向量 $(x_0, x_1)$ 视为复数 $z = x_0 + i x_1$。
根据欧拉公式，旋转角度 $\theta$ 相当于乘以 $e^{i \theta}$：

\[
z_{\text{rot}} = z \cdot e^{i \theta} = (x_0 + i x_1)(\cos\theta + i\sin\theta) = (x_0\cos\theta - x_1\sin\theta) + i(x_0\sin\theta + x_1\cos\theta)
\]

转换为实数矩阵乘法形式：

\[
R_{\theta} = \begin{bmatrix}
\cos\theta & -\sin\theta \\
\sin\theta & \cos\theta
\end{bmatrix}
\]

在位置 $m$ 处，旋转角度为 $m \theta$：

\[
q_m = R_{m \theta} q
\]

在位置 $n$ 处，旋转角度为 $n \theta$：

\[
k_n = R_{n \theta} k
\]

\subsection{为什么内积只取决于相对距离？}

让我们计算旋转后的内积：

\[
\langle q_m, k_n \rangle = (R_{m \theta} q)^T (R_{n \theta} k) = q^T (R_{m \theta}^T R_{n \theta}) k
\]

注意旋转矩阵的正交性质：逆时针旋转 $m\theta$ 的转置矩阵，就等于顺时针旋转，即 $R_{m \theta}^T = R_{-m \theta}$！
根据旋转矩阵的相乘性质（先转 $-m\theta$ 再转 $n\theta$）：

\[
R_{m \theta}^T R_{n \theta} = R_{-m \theta} R_{n \theta} = R_{(n - m) \theta}
\]

代入内积式：

\[
\langle q_m, k_n \rangle = q^T R_{(n - m) \theta} k = g(q, k, n - m) \quad \blacksquare
\]

\textbf{数学证明完毕！}
通过在二维平面上纯粹旋转一个角度 $m\theta$：
\begin{enumerate}
  \item 向量本身的长度（模长）没有被拉伸或破坏；
  \item 两个词点积时，两者的绝对位置 $m$ 和 $n$ 自动相消，\textbf{唯独留下了相对距离 $n - m$}！
  \item 当高维空间拆成许多个 2D 坐标对，每个对使用不同的旋转频率 $\theta_i = 10000^{-2(i-1)/d}$ 时，会出现一种\textbf{有条件的}"长程衰减"——这一点经常被说错，要精确表述：
\end{enumerate}

\begin{itemize}
  \item 把第 $i$ 对分量的贡献写成 $A_i \cos(\Delta\theta_i) + B_i \sin(\Delta\theta_i)$（$\Delta$ 是相对距离）。如果 $q$ 与 $k$ \textbf{是相关的}（语义上"匹配"，$A_i$ 平均为正），那么点积的\textbf{期望}是一堆不同频率余弦的正权重叠加，$\Delta$ 变大后各频率相位错开，\textbf{期望值随距离衰减}：同样相似的两个词，离得越远得到的注意力加成越小。
  \item 但如果 $q$ 与 $k$ \textbf{不相关}，点积的分布与距离\textbf{完全无关}（旋转不改变各向同性随机向量的分布），没有任何衰减。对某一对具体的向量，点积随距离剧烈震荡，远处的值完全可能比近处更大。
  \item RoFormer 论文 §3.4.3 证明的是点积的一个\textbf{上界}随距离衰减，并不是点积本身一定变小。
  \item 所以 RoPE 的衰减是"对匹配信号的软性偏置"，不是硬性的距离惩罚。长上下文扩展（PI / NTK / YaRN，第 5 篇 §5.4）要解决的是另一个问题：超过训练长度后，低频维度会转到训练时从未见过的角度（分布外）。
\end{itemize}

\vspace{0.8em}\hrule\vspace{0.8em}

\section{[目标] 总结与下一章预告}

在本篇中，我们逐一完成了 4 大核心公式的严格数学推导：
\begin{enumerate}
  \item \textbf{$Q, K, V$ 拆分}：打破对称性，建立检索信息流；
  \item \textbf{除以 $\sqrt{d_k}$}：证明了独立变量点积方差随维度线性增长，通过平方根归一化彻底拯救了 Softmax 的梯度消失；
  \item \textbf{因果掩码 $-\infty$}：利用 $e^{-\infty}=0$ 斩断未来信息泄露；
  \item \textbf{残差连接}：通过常数导数高速公路根治深度退化；
  \item \textbf{RoPE 旋转矩阵}：通过二维复平面旋转证明了相对位置内积恒等性。
\end{enumerate}

至此，Transformer 的核心公式你已经推导了一遍！

\begin{tipBox}[核心直觉与思考]
{}[重点] \textbf{但注意}：本篇讲的是 2017 年原版论文的公式。现代大模型（Llama、Qwen、DeepSeek）在此基础上做了许多改动——Pre-LN、RMSNorm、SwiGLU、GQA、MLA、MoE——每一处都有明确的历史动机，见 \textbf{第 5 篇}。
想亲手验证本篇的 RoPE 性质（包括"衰减只发生在 q、k 相关时的期望值上，不相关时完全不衰减"这一容易被误解的点），运行 \texttt{lab04\_transformer\_microscope/02\_rope\_and\_attention\_microscope.py}。
\end{tipBox}

在下一篇（第 5 篇）中，我们先看这些公式在 2017 年之后被改了哪些地方、为什么改；再下一篇（第 6 篇）跨越到实际运行：\textbf{到底什么是大模型推理（Inference）？为什么它在首字（Prefill）和后续生成（Decode）表现出完全不同的运行规律？KV Cache 在代数上是如何推演出来的？}

\vspace{0.8em}\hrule\vspace{0.8em}

\textbf{上一篇}：\textbf{第 3 篇：神经网络如何学习：导数、链式法则、反向传播与交叉熵} ｜ \textbf{下一篇}：\textbf{第 5 篇：从原版 Transformer 到 Llama / DeepSeek：每一处改动的历史动机} ｜ \textbf{总路线}：\textbf{LEARNING\_PATH.md}
